\documentclass[12pt,a4paper]{report}

\usepackage{fontspec}
\usepackage{polyglossia}
\setdefaultlanguage{romanian}
\usepackage[margin=2.25cm]{geometry}
\usepackage{setspace}
\usepackage{amsmath}
\usepackage{amssymb}
\usepackage{enumitem}
\usepackage{longtable}
\usepackage{array}
\usepackage{booktabs}
\usepackage{xcolor}
\usepackage{listings}
\usepackage[
    colorlinks,
    linkcolor={black},
    citecolor={black},
    urlcolor={blue}
]{hyperref}

\onehalfspacing
\setlength{\emergencystretch}{3em}
\setlength{\parskip}{0.15em}
\setlist{nosep,leftmargin=*}

\definecolor{codebackground}{rgb}{0.96,0.96,0.96}
\definecolor{codecomment}{rgb}{0.23,0.45,0.30}
\definecolor{codekeyword}{rgb}{0.10,0.20,0.55}
\definecolor{codestring}{rgb}{0.55,0.15,0.15}

\lstdefinestyle{codestyle}{
    backgroundcolor=\color{codebackground},
    basicstyle=\ttfamily\scriptsize,
    breaklines=true,
    frame=single,
    columns=fullflexible,
    keepspaces=true,
    showstringspaces=false,
    numbers=left,
    numberstyle=\tiny\color{gray},
    keywordstyle=\color{codekeyword}\bfseries,
    commentstyle=\color{codecomment},
    stringstyle=\color{codestring},
    tabsize=4
}
\lstset{style=codestyle}

\newcommand{\codechunk}[2]{%
\subsection*{Liniile #1--#2}
\lstinputlisting[language=Python,firstline=#1,lastline=#2,firstnumber=#1]{../src/dsp/stft_utils.py}
}

\title{Explicarea fisierului \texttt{stft\_utils.py}}
\author{}
\date{}

\begin{document}

\maketitle
\tableofcontents
\clearpage

\chapter{Rolul fisierului \texttt{stft\_utils.py}}

Fisierul \path{src/dsp/stft_utils.py} contine functiile fundamentale pentru
trecerea dintre domeniul timp si domeniul timp-frecventa. Aceste functii sunt
folosite in mai multe zone ale proiectului: metodele clasice DSP, inferenta
MaskNet, training-ul cu masti spectrale si varianta Complex MaskNet.

STFT-ul transforma semnalul audio \(x[n]\) intr-o reprezentare complexa
\(X[k,m]\), unde \(k\) este frecventa si \(m\) este cadrul temporal:

\[
X[k,m] =
\sum_{n=0}^{L-1} x[n+mH]w[n]e^{-j2\pi kn/N}.
\]

ISTFT-ul reconstruieste un semnal temporal din spectrograma complexa, prin
suprapunerea si adunarea cadrelor:

\[
\hat{x}[n] = \operatorname{ISTFT}(X[k,m]).
\]

Fisierul contine doua familii de functii:

\begin{itemize}
    \item implementari NumPy simple pentru STFT/ISTFT si operatii magnitudine-faza;
    \item implementari PyTorch folosite in modelul neuronal si inferenta GPU/CPU.
\end{itemize}

\chapter{Importuri}

\codechunk{1}{3}

\begin{description}
    \item[Linia 1.] Importa NumPy, folosit pentru implementarea clasica STFT/ISTFT
    si pentru operatii pe vectori audio in format \texttt{np.ndarray}.

    \item[Linia 2.] Importa PyTorch, necesar pentru functiile STFT/ISTFT folosite
    in inferenta MaskNet si in training.

    \item[Linia 3.] Importa \texttt{config}, de unde sunt preluati parametrii
    impliciti: \texttt{FRAME\_LEN}, \texttt{HOP\_LEN} si \texttt{N\_FFT}.
\end{description}

\chapter{Segmentarea semnalului}

\codechunk{5}{25}

\begin{description}
    \item[Linia 5.] \texttt{frame\_signal} imparte un semnal temporal in cadre
    suprapuse. Aceasta etapa este baza STFT-ului.

    \item[Liniile 6--9.] Daca nu sunt furnizate valori explicite pentru lungimea
    cadrului si hop, sunt folosite valorile din configuratia proiectului.

    \item[Liniile 10--12.] Daca semnalul are mai multe canale, este convertit la
    mono prin medierea canalelor. Apoi este convertit la \texttt{float32}, format
    eficient si suficient pentru procesare audio.

    \item[Liniile 13--17.] Daca semnalul este mai scurt decat un cadru, este completat
    cu zerouri pana ajunge la lungimea minima necesara.

    \item[Liniile 18--20.] Calculeaza numarul de cadre:
    \[
    M = 1 + \left\lfloor \frac{N-L}{H} \right\rfloor.
    \]
    Apoi aloca matricea \texttt{frames} cu forma \((M,L)\).

    \item[Liniile 21--24.] Parcurge cadrele si copiaza esantioanele corespunzatoare
    fiecarui segment.

    \item[Linia 25.] Intoarce matricea de cadre, pregatita pentru aplicarea ferestrei
    si FFT.
\end{description}

\chapter{STFT si ISTFT in NumPy}

\codechunk{27}{38}

\begin{description}
    \item[Linia 27.] \texttt{stft} calculeaza Short-Time Fourier Transform pentru
    un semnal NumPy.

    \item[Liniile 28--33.] Preia valorile implicite pentru \texttt{n\_fft},
    \texttt{hop\_length} si \texttt{win\_length} din \texttt{config.py}.

    \item[Linia 34.] Imparte semnalul in cadre folosind \texttt{frame\_signal}.

    \item[Liniile 35--36.] Creeaza o fereastra Hanning si o aplica fiecarui cadru.
    Fereastra reduce discontinuitatile de la marginile cadrelor.

    \item[Linia 37.] Aplica FFT real, \texttt{np.fft.rfft}, pe fiecare cadru. Pentru
    un FFT de dimensiune \(N\), rezultatul are \(N/2+1\) bin-uri.

    \item[Linia 38.] Intoarce spectrograma transpusă, cu forma \((F,T)\), adica
    frecventa pe prima axa si timpul pe a doua.
\end{description}

\codechunk{40}{70}

\begin{description}
    \item[Linia 40.] \texttt{istft} reconstruieste un semnal temporal dintr-o
    spectrograma complexa NumPy.

    \item[Liniile 41--46.] Completeaza parametrii lipsa cu valorile implicite din
    configuratie.

    \item[Liniile 47--52.] Transpune spectrograma pentru a avea cadrele pe prima
    axa, aplica inverse FFT real si pastreaza doar \texttt{win\_length} esantioane
    din fiecare cadru.

    \item[Linia 53.] Reaplica fereastra Hanning cadrelor reconstruite.

    \item[Liniile 54--61.] Initializeaza semnalul de iesire si face overlap-add:
    fiecare cadru este adunat in pozitia lui temporala. In paralel, \texttt{win\_sum}
    acumuleaza ferestrele pentru normalizare.

    \item[Liniile 62--63.] Normalizeaza semnalul in zonele unde suma ferestrelor
    este nenula. Aceasta compenseaza suprapunerea cadrelor.

    \item[Liniile 64--69.] Daca este ceruta o lungime exacta, semnalul este taiat
    sau completat cu zerouri.

    \item[Linia 70.] Intoarce waveform-ul reconstruit.
\end{description}

\chapter{Operatii magnitudine-faza}

\codechunk{72}{78}

\begin{description}
    \item[Liniile 72--75.] \texttt{mag\_phase} separa o spectrograma complexa in
    magnitudine si faza:
    \[
    |X[k,m]| = \sqrt{\Re(X)^2 + \Im(X)^2},
    \quad
    \angle X[k,m] = \arg(X[k,m]).
    \]

    \item[Liniile 77--78.] \texttt{from\_mag\_phase} reconstruieste spectrograma
    complexa din magnitudine si faza:
    \[
    X[k,m] = |X[k,m]|e^{j\angle X[k,m]}.
    \]
\end{description}

\chapter{Conversii pentru spectrograme complexe}

\codechunk{81}{100}

\begin{description}
    \item[Liniile 81--89.] \texttt{complex\_to\_channels} converteste un tensor
    complex PyTorch intr-un tensor cu doua canale: partea reala si partea imaginara.
    Aceasta forma este necesara pentru reteaua convolutionala, deoarece layer-ele
    \texttt{Conv2d} lucreaza cu canale reale, nu direct cu numere complexe.

    \item[Forme acceptate.] Pentru o spectrograma \((F,T)\), rezultatul este
    \((2,F,T)\). Pentru un batch \((B,F,T)\), rezultatul este \((B,2,F,T)\).

    \item[Liniile 92--100.] \texttt{channels\_to\_complex} face operatia inversa:
    transforma doua canale reale intr-un tensor complex. Daca forma nu este cea
    asteptata, functia ridica o eroare explicita.
\end{description}

\chapter{Complex Ratio Mask}

\codechunk{103}{132}

\begin{description}
    \item[Liniile 103--112.] \texttt{compute\_complex\_ratio\_mask} calculeaza masca
    complexa care transforma STFT-ul zgomotos in STFT-ul curat. Iesirea are doua
    canale: masca reala si masca imaginara.

    \item[Liniile 114--117.] Extrage partile reale si imaginare ale STFT-ului
    zgomotos si ale STFT-ului curat.

    \item[Liniile 119--121.] Calculeaza formula CRM:
    \[
    M_r = \frac{S_rY_r + S_iY_i}{Y_r^2 + Y_i^2 + \epsilon},
    \qquad
    M_i = \frac{S_iY_r - S_rY_i}{Y_r^2 + Y_i^2 + \epsilon}.
    \]
    Unde \(S\) este STFT-ul curat, iar \(Y\) este STFT-ul zgomotos.

    \item[Liniile 123--127.] Alege dimensiunea pe care se stivuiesc cele doua
    canale si limiteaza valorile mastii intre \(-\texttt{clip\_value}\) si
    \(\texttt{clip\_value}\). Limitarea ajuta la stabilitatea training-ului.

    \item[Liniile 130--132.] \texttt{apply\_complex\_mask} transforma masca din doua
    canale in tensor complex si o inmulteste cu spectrograma zgomotoasa:
    \[
    \hat{S}[k,m] = M_{\mathbb{C}}[k,m]Y[k,m].
    \]
\end{description}

\chapter{STFT in PyTorch}

\codechunk{135}{194}

\begin{description}
    \item[Liniile 135--155.] \texttt{compute\_stft} este varianta PyTorch a STFT-ului.
    Poate intoarce fie magnitudine si faza, fie spectrograma complexa.

    \item[Liniile 156--161.] Completeaza parametrii lipsa cu valorile din
    \texttt{config.py}.

    \item[Liniile 163--167.] Daca semnalul are forma \((T)\), adauga dimensiune de
    batch si retine ca iesirea trebuie ulterior stransa inapoi.

    \item[Linia 169.] Creeaza fereastra Hann direct pe dispozitivul tensorului,
    CPU sau GPU.

    \item[Liniile 171--180.] Apeleaza \texttt{torch.stft}. Parametrul
    \texttt{return\_complex=True} produce tensor complex, \texttt{center=True}
    aplica padding in jurul semnalului, iar \texttt{pad\_mode='reflect'} reduce
    artefactele de margine.

    \item[Liniile 182--185.] Daca se cere spectrograma complexa, functia o intoarce
    direct, eliminand dimensiunea de batch daca intrarea a fost 1D.

    \item[Liniile 187--194.] In cazul standard, separa magnitudinea si faza si le
    intoarce cu aceeasi conventie dimensionala ca intrarea.
\end{description}

\chapter{ISTFT din magnitudine si faza}

\codechunk{197}{250}

\begin{description}
    \item[Liniile 197--218.] \texttt{inverse\_stft} reconstruieste waveform-ul din
    magnitudine si faza PyTorch. Este folosita de modelele care prezic doar masca
    de magnitudine.

    \item[Liniile 219--224.] Completeaza parametrii lipsa cu valorile implicite.

    \item[Liniile 226--231.] Daca intrarea nu are batch, adauga batch temporar si
    retine ca trebuie eliminat la final.

    \item[Linia 233.] Reconstruieste spectrograma complexa:
    \[
    X[k,m] = |X[k,m]|e^{j\phi[k,m]}.
    \]

    \item[Liniile 235--245.] Creeaza fereastra Hann si apeleaza \texttt{torch.istft}.
    Parametrul \texttt{length} permite obtinerea unui waveform cu lungime exacta.

    \item[Liniile 247--250.] Elimina dimensiunea de batch daca a fost adaugata
    artificial si intoarce semnalul reconstruit.
\end{description}

\chapter{ISTFT complex}

\codechunk{253}{292}

\begin{description}
    \item[Liniile 253--265.] \texttt{inverse\_complex\_stft} reconstruieste semnalul
    direct din spectrograma complexa. Este folosita in varianta Complex MaskNet.

    \item[Liniile 266--271.] Completeaza parametrii STFT lipsa cu valorile implicite.

    \item[Liniile 273--277.] Adauga dimensiune de batch daca spectrograma are doar
    forma \((F,T)\).

    \item[Liniile 279--288.] Creeaza fereastra Hann si apeleaza \texttt{torch.istft}
    direct pe tensorul complex.

    \item[Liniile 290--292.] Elimina dimensiunea de batch daca este cazul si intoarce
    waveform-ul final.
\end{description}

\chapter{Concluzie}

\texttt{stft\_utils.py} este unul dintre cele mai importante fisiere DSP ale
proiectului. El ofera reprezentarea timp-frecventa folosita atat de metodele
clasice, cat si de MaskNet. Functiile NumPy sunt utile pentru metode DSP si
analiza, iar functiile PyTorch sunt esentiale pentru training si inferenta.

Prin includerea functiilor pentru CRM, fisierul sustine si varianta Complex
MaskNet, unde modelul nu mai modifica doar magnitudinea, ci poate estima corectii
complexe asupra spectrogramei. Astfel, acest modul este puntea matematica dintre
semnalul audio brut si reprezentarile pe care modelul le poate invata.

\end{document}
